\documentclass[UTF8,fontset=fandol,11pt]{ctexart}
\usepackage[a4paper,margin=24mm]{geometry}
\usepackage{amsmath,amssymb,amsthm,graphicx,tabularx,booktabs,hyperref}
\newtheorem{definition}{定义}[section]
\newtheorem{theorem}[definition]{定理}
\newtheorem{example}[definition]{例题}
\newtheorem{exercise}[definition]{练习}
\hypersetup{unicode=true,colorlinks=true,linkcolor=blue}
\title{最小二乘法讲义}
\date{}
\begin{document}
\maketitle
\section{误差如何度量：从残差到常数拟合}\label{sec:S01}
% UNIT S01.U01
对同一个量重复测量时，读数往往并不完全相同。如果要用一个数代表这些读数，应当如何选择？最小二乘法先规定误差的度量，再寻找使这种误差最小的模型参数。本节从常数拟合入手，只使用有限求和、一元求导和配方法。\par
% UNIT S01.U02
\begin{definition}
设 $n\geq 1$ 为整数，$y_i\in\mathbb{R}$ 是第 $i$ 次观测的响应值，$i=1,\ldots,n$。若记录中还有自变量，则用 $x_i\in\mathbb{R}$ 表示第 $i$ 次观测的输入。给定模型及其参数后，模型对第 $i$ 次观测给出的数值记为 $\hat{y}_i$，称为拟合值。

约定残差为观测值减去拟合值：
\[
r_i=y_i-\hat{y}_i,\qquad i=1,\ldots,n.
\]
因此，$r_i>0$ 表示模型低估了该观测值，$r_i<0$ 表示模型高估了该观测值。残差平方和定义为
\[
Q=\sum_{i=1}^{n}r_i^2=\sum_{i=1}^{n}(y_i-\hat{y}_i)^2.
\]
最小二乘准则是：在选定模型的允许参数范围内，寻找使 $Q$ 取得最小值的参数。这里最小化的是所选模型对已有数据的残差平方和，而不是保证每个残差都分别最小。
\end{definition}
% UNIT S01.U03
为什么不直接用残差之和度量拟合误差？正负残差可能互相抵消。例如，两个残差为 $1$ 与 $-1$ 时，残差之和为零，但两个拟合值都没有命中观测值。平方后各项非负，不能发生这种抵消；并且
\[
\sum_{i=1}^{n}r_i^2=0\quad\Longleftrightarrow\quad r_i=0\ \text{对所有 }i=1,\ldots,n.
\]
平方也会使较大的残差更显著地影响目标函数。如果响应值有计量单位，残差与响应值同单位，而残差平方和具有该单位的平方。比较不同候选模型的平方和时，本节始终使用同一组数据和相同单位。\par
% UNIT S01.U04
\begin{definition}
设 $n\geq1$ 为整数，观测值为 $y_1,\ldots,y_n\in\mathbb{R}$。常数模型用同一个待定实数 $c$ 拟合所有观测值，即
\[
\hat{y}_i=c,\qquad i=1,\ldots,n,\qquad c\in\mathbb{R}.
\]
相应的目标函数是
\[
Q(c)=\sum_{i=1}^{n}(y_i-c)^2,\qquad c\in\mathbb{R}.
\]
这些观测值的样本算术平均数记为
\[
\bar{y}=\frac{1}{n}\sum_{i=1}^{n}y_i.
\]
在优化过程中，观测值以及由它们算出的 $\bar{y}$ 都是固定的，只有参数 $c$ 发生变化。
\end{definition}
% UNIT S01.U05
\begin{theorem}
设 $n\geq1$ 为整数，$y_1,\ldots,y_n\in\mathbb{R}$，并令 $\bar{y}=\frac{1}{n}\sum_{i=1}^{n}y_i$。则常数模型的残差平方和满足：对每个 $c\in\mathbb{R}$，
\[
Q(c)=\sum_{i=1}^{n}(y_i-\bar{y})^2+n(c-\bar{y})^2.
\]
因此，$Q(c)$ 在且仅在 $c=\bar{y}$ 时取得全局最小值，最小值为
\[
Q(\bar{y})=\sum_{i=1}^{n}(y_i-\bar{y})^2.
\]
\end{theorem}
% UNIT S01.U06
\begin{proof}
先用一元求导寻找候选参数。对任意 $c\in\mathbb{R}$，有限项逐项求导得
\[
Q'(c)=2\sum_{i=1}^{n}(c-y_i)=2n(c-\bar{y}).
\]
所以 $Q'(c)=0$ 当且仅当 $c=\bar{y}$。由于参数范围是整个实数轴，可微函数若在某点取得最小值，该点必须满足导数为零。不过，求出驻点本身还不是全局最小性的证明。

下面直接比较所有参数对应的函数值。由平均数的定义，
\[
\sum_{i=1}^{n}(y_i-\bar{y})=\sum_{i=1}^{n}y_i-n\bar{y}=0.
\]
于是，对任意 $c\in\mathbb{R}$，平方展开得到
\[
\begin{aligned}
Q(c)&=\sum_{i=1}^{n}\bigl((y_i-\bar{y})+(\bar{y}-c)\bigr)^2\\
&=\sum_{i=1}^{n}(y_i-\bar{y})^2
 +2(\bar{y}-c)\sum_{i=1}^{n}(y_i-\bar{y})
 +n(\bar{y}-c)^2\\
&=Q(\bar{y})+n(c-\bar{y})^2.
\end{aligned}
\]
因为 $n>0$，最后一项对所有实数 $c$ 都非负，且仅在 $c=\bar{y}$ 时为零。这既证明了最小值存在，也证明了全局最小点唯一。
\end{proof}
% UNIT S01.U07
\begin{example}
某个保持不变的量经过三次测量，读数依次为 $2,3,7$，单位相同。用常数模型拟合这些读数：写出目标函数，求最优常数及最小残差平方和，并与候选常数 $c=3$ 比较。
\end{example}
% UNIT S01.U08
\par\noindent\textbf{解：}对 $c\in\mathbb{R}$，目标函数为
\[
\begin{aligned}
Q(c)&=(2-c)^2+(3-c)^2+(7-c)^2\\
&=3c^2-24c+62\\
&=3(c-4)^2+14.
\end{aligned}
\]
三个读数的平均数是 $\bar{y}=(2+3+7)/3=4$，所以唯一最优常数为 $c=4$。这也可从 $Q'(c)=6c-24$ 和上述配方看出。

将两个候选常数的残差列在一起：
\[
\begin{array}{c|rrr|r}
\text{候选常数}&r_1&r_2&r_3&Q(c)\\ \hline
c=3&-1&0&4&17\\
c=4&-2&-1&3&14
\end{array}
\]
因此最小平方和为 $Q(4)=4+1+9=14$，而 $Q(3)=1+0+16=17$。从 $c=3$ 改为 $c=4$ 后，前两个残差的绝对值反而增大，但总平方和减小了 $3$。这说明最小二乘法优化的是整体指标，并不要求每次测量的拟合误差都同时减小。\par
% UNIT S01.U09
常数拟合给出了一个可重复使用的思路：先明确参数与目标函数，再求导寻找候选点，最后用平方和分解证明全局最优性。恒等式 $Q(c)-Q(\bar{y})=n(c-\bar{y})^2$ 还定量说明了偏离最优参数的代价。

最优常数模型的残差之和为 $\sum_{i=1}^{n}(y_i-\bar{y})=0$，但残差平方和不一定为零；它为零当且仅当所有观测值相同。第\ref{sec:S02}节将把这里的常数拟合结论用于直线拟合，仍然只需要一元求导与配方。\par
% UNIT S01.U10
\begin{exercise}
\label{mathagent:exercise:79e91357ee05b8b9aa4c657d553157f1ef411624d282c3f2337dd0e01453ce2d}
对同单位的四次测量值 $1,2,4,5$ 作常数拟合。写出 $Q(c)$ 的平方和形式，求最优常数、各个残差及最小平方和；再计算 $Q(2)$，并用平方和分解解释它与最小值之差。本题答案统一置于讲义末尾。
\end{exercise}
% UNIT S01.U11
\begin{exercise}
\label{mathagent:exercise:9e7dbf7ef17df1506f236bbf1ad6ff8c14ff13065df89e12b796d18475228be7}
设 $n\geq1$ 为整数，$y_1,\ldots,y_n\in\mathbb{R}$，常数模型的残差为 $r_i=y_i-c$，其中 $c\in\mathbb{R}$。判断并说明理由：

（1）若 $\sum_{i=1}^{n}r_i=0$，则 $c$ 是该常数模型的唯一最优参数。

（2）若 $\sum_{i=1}^{n}r_i=0$，则每一个残差都等于零。

（3）当 $n=1$ 时，常数模型仍有唯一最优参数。

本题答案统一置于讲义末尾。
\end{exercise}
\section{直线最小二乘：消去截距并求解正规方程}\label{sec:S02}
% UNIT S02.U01
第\ref{sec:S01}节用一个常数代表一组观测值。本节允许拟合值随输入作线性变化，用直线描述数据的整体趋势。虽然有截距和斜率两个参数，我们仍只用一元微积分：先固定斜率求最佳截距，再求最佳斜率，最后用配方证明全局最优性。\par
% UNIT S02.U02
\begin{definition}
设 $n\geq1$ 为整数，给定数据点 $(x_i,y_i)\in\mathbb{R}^2$，$i=1,\ldots,n$。带截距的直线拟合模型为
\[
\hat{y}_i=a+bx_i,\qquad a,b\in\mathbb{R},
\]
其中 $a$ 是截距，$b$ 是斜率。残差与目标函数分别为
\[
r_i=y_i-a-bx_i,\qquad
Q(a,b)=\sum_{i=1}^{n}(y_i-a-bx_i)^2.
\]
最小二乘直线指使 $Q(a,b)$ 在所有实参数对 $(a,b)$ 中取得最小值的直线。这里截距也是待定参数，并未要求直线经过原点。
\end{definition}
% UNIT S02.U03
\begin{definition}
对上述数据，定义均值
\[
\bar{x}=\frac1n\sum_{i=1}^{n}x_i,\qquad
\bar{y}=\frac1n\sum_{i=1}^{n}y_i.
\]
中心化数据是 $x_i-\bar{x}$ 与 $y_i-\bar{y}$，其各自的和均为零。定义
\[
S_{xx}=\sum_{i=1}^{n}(x_i-\bar{x})^2,\qquad
S_{xy}=\sum_{i=1}^{n}(x_i-\bar{x})(y_i-\bar{y}).
\]
平方展开并利用均值定义，可得便于手算的形式
\[
S_{xx}=\sum_{i=1}^{n}x_i^2-n\bar{x}^2,\qquad
S_{xy}=\sum_{i=1}^{n}x_i y_i-n\bar{x}\bar{y}.
\]
这些量都只由观测数据决定，与待定参数 $a,b$ 无关。
\end{definition}
% UNIT S02.U04
固定任意 $b\in\mathbb{R}$，把 $y_i-bx_i$ 看作新的观测值。它们的平均数是 $\bar{y}-b\bar{x}$。由第\ref{sec:S01}节的常数拟合结论，此时唯一最优截距是
\[
a=\bar{y}-b\bar{x}.
\]
更具体地，对所有 $a,b\in\mathbb{R}$，有分解
\[
Q(a,b)=F(b)+n(a+b\bar{x}-\bar{y})^2,
\]
其中
\[
\begin{aligned}
F(b)&=Q(\bar{y}-b\bar{x},b)\\
&=\sum_{i=1}^{n}\bigl((y_i-\bar{y})-b(x_i-\bar{x})\bigr)^2\\
&=\sum_{i=1}^{n}(y_i-\bar{y})^2-2bS_{xy}+b^2S_{xx}.
\end{aligned}
\]
因此，先把截距选为对应斜率下的最优值，再最小化一元函数 $F$，不会漏掉更好的参数对。若 $S_{xx}>0$，则
\[
F'(b)=2S_{xx}b-2S_{xy},
\]
导数为零给出候选斜率 $b=S_{xy}/S_{xx}$。下面用配方说明它确实给出唯一全局最优解。\par
% UNIT S02.U05
\begin{definition}
对给定的 $n$ 个实数据点，本模型关于未知参数 $a,b\in\mathbb{R}$ 的正规方程是二元一次方程组
\[
\begin{cases}
na+b\displaystyle\sum_{i=1}^{n}x_i=\displaystyle\sum_{i=1}^{n}y_i,\\
a\displaystyle\sum_{i=1}^{n}x_i+b\displaystyle\sum_{i=1}^{n}x_i^2=\displaystyle\sum_{i=1}^{n}x_i y_i.
\end{cases}
\]
按残差 $r_i=y_i-a-bx_i$ 的约定，这也就是
\[
\sum_{i=1}^{n}r_i=0,\qquad \sum_{i=1}^{n}x_i r_i=0.
\]
这两个等式不表示每个残差分别为零。它们与最小二乘解的联系将在下面通过消元与配方建立，不需要偏导数。
\end{definition}
% UNIT S02.U06
\begin{theorem}
设 $n\geq1$ 为整数，$(x_i,y_i)\in\mathbb{R}^2$，$i=1,\ldots,n$。条件 $S_{xx}>0$ 等价于 $x_1,\ldots,x_n$ 不全相同。在此条件下，$Q(a,b)$ 在 $\mathbb{R}^2$ 上有唯一全局最小点
\[
b_* =\frac{S_{xy}}{S_{xx}},\qquad a_* =\bar{y}-b_*\bar{x}.
\]
参数对 $(a_*,b_*)$ 也是正规方程的唯一解。对任意 $a,b\in\mathbb{R}$，有
\[
Q(a,b)=Q(a_*,b_*)+S_{xx}(b-b_*)^2+n(a+b\bar{x}-\bar{y})^2,
\]
其中
\[
Q(a_*,b_*)=\sum_{i=1}^{n}(y_i-\bar{y})^2-\frac{S_{xy}^2}{S_{xx}}.
\]
\end{theorem}
% UNIT S02.U07
\begin{proof}
由于 $S_{xx}$ 是非负项之和，$S_{xx}=0$ 当且仅当每个 $x_i-\bar{x}=0$，也就是所有 $x_i$ 相同。因此 $S_{xx}>0$ 当且仅当自变量观测值不全相同。

以下设 $S_{xx}>0$。对消去截距后的一元二次函数配方，得到对所有实数 $b$ 成立的恒等式
\[
\begin{aligned}
F(b)
&=\sum_{i=1}^{n}(y_i-\bar{y})^2-2bS_{xy}+b^2S_{xx}\\
&=\sum_{i=1}^{n}(y_i-\bar{y})^2-\frac{S_{xy}^2}{S_{xx}}
  +S_{xx}\left(b-\frac{S_{xy}}{S_{xx}}\right)^2.
\end{aligned}
\]
令 $b_*=S_{xy}/S_{xx}$、$a_*=\bar{y}-b_*\bar{x}$，再代入前面的截距分解，便得到定理中的 $Q(a,b)$ 分解式。右侧两个附加平方项都非负；它们同时为零当且仅当 $b=b_*$ 且 $a+b\bar{x}-\bar{y}=0$，即 $(a,b)=(a_*,b_*)$。故最小值存在，且全局最小点唯一。

最后用消元联系正规方程。第一条正规方程除以 $n>0$ 后等价于
\[
a+b\bar{x}=\bar{y}.
\]
将第二条正规方程减去第一条的 $\bar{x}$ 倍，利用 $\sum x_i=n\bar{x}$，得到
\[
b\left(\sum_{i=1}^{n}x_i^2-n\bar{x}^2\right)
=\sum_{i=1}^{n}x_i y_i-n\bar{x}\bar{y},
\]
也就是 $bS_{xx}=S_{xy}$。上述消元可逆，所以正规方程等价于
\[
a+b\bar{x}=\bar{y},\qquad bS_{xx}=S_{xy}.
\]
由于 $S_{xx}>0$，其唯一解正是 $(a_*,b_*)$。这同时证明了最优参数满足正规方程，以及在本条件下解正规方程足以求得唯一最优参数。
\end{proof}
% UNIT S02.U08
\begin{example}
对数据点 $(0,1)$、$(1,2)$、$(2,2)$ 求带截距的最小二乘直线。列出均值、中心化求和、最优参数、各个残差及最小残差平方和，并核对两条正规方程。
\end{example}
% UNIT S02.U09
\par\noindent\textbf{解：}均值为
\[
\bar{x}=\frac{0+1+2}{3}=1,\qquad
\bar{y}=\frac{1+2+2}{3}=\frac53.
\]
中心化计算如下：
\[
\begin{array}{c|r|r|r|r}
i&x_i-\bar{x}&y_i-\bar{y}&(x_i-\bar{x})^2&(x_i-\bar{x})(y_i-\bar{y})\\ \hline
1&-1&-2/3&1&2/3\\
2&0&1/3&0&0\\
3&1&1/3&1&1/3
\end{array}
\]
因此 $S_{xx}=2>0$、$S_{xy}=1$，唯一最优参数为
\[
b_* =\frac12,\qquad a_* =\frac53-\frac12=\frac76.
\]
最小二乘直线是 $\hat{y}=\frac76+\frac12x$。逐点计算拟合值和残差：
\[
\begin{array}{c|r|r|r|r}
i&x_i&y_i&\hat{y}_i&r_i\\ \hline
1&0&1&7/6&-1/6\\
2&1&2&5/3&1/3\\
3&2&2&13/6&-1/6
\end{array}
\]
故最小平方和为
\[
Q(a_*,b_*)=\frac1{36}+\frac19+\frac1{36}=\frac16.
\]
原始数据给出 $\sum x_i=3$、$\sum y_i=5$、$\sum x_i^2=5$、$\sum x_i y_i=6$，所以正规方程为
\[
3a+3b=5,\qquad 3a+5b=6.
\]
代入最优参数，分别有 $7/2+3/2=5$ 与 $7/2+5/2=6$。等价地，残差满足
\[
-\frac16+\frac13-\frac16=0,\qquad
0\left(-\frac16\right)+1\left(\frac13\right)+2\left(-\frac16\right)=0.
\]
还可以把全局最优性具体写为：对所有 $a,b\in\mathbb{R}$，
\[
\begin{aligned}
Q(a,b)&=(1-a)^2+(2-a-b)^2+(2-a-2b)^2\\
&=\frac16+2\left(b-\frac12\right)^2+3\left(a+b-\frac53\right)^2.
\end{aligned}
\]
本例中三个点不在同一条直线上，因而最优残差也不全为零；最小二乘拟合并不是强行让直线通过每一个数据点。\par
% UNIT S02.U10
手算时可按以下顺序进行：先算 $\bar{x},\bar{y}$，再算 $S_{xx},S_{xy}$；确认 $S_{xx}>0$ 后才作除法求 $b_*$，随后求 $a_*$；最后计算残差平方和，并用正规方程核对参数。

注意，$S_{xy}=0$ 并不意味着公式失效：只要 $S_{xx}>0$，最优斜率就是零，最优直线为水平线 $\hat{y}=\bar{y}$。真正不能直接使用除法公式的是 $S_{xx}=0$ 的情形。它的参数唯一性问题，以及残差和均值点的解释，将在第\ref{sec:S03}节讨论。\par
% UNIT S02.U11
\begin{exercise}
\label{mathagent:exercise:332ae2e99c79cf97370ba4732967ff23f2a2c4a4b36b98eda8e910755bf43b8b}
对数据点 $(0,1)$、$(1,1)$、$(2,3)$ 求带截距的最小二乘直线。计算 $\bar{x},\bar{y},S_{xx},S_{xy}$，列出残差及最小平方和，并核对正规方程。答案统一置于讲义末尾。
\end{exercise}
% UNIT S02.U12
\begin{exercise}
\label{mathagent:exercise:b4c9c9a1fb82bcef052fdbf6e1539831ed88cee016dbcbf83a50e5e8d2311be0}
设 $n\geq2$ 为整数，给定实数据点 $(x_i,y_i)$，$i=1,\ldots,n$，且 $S_{xx}>0$。判断下列说法并说明理由。

（1）若 $S_{xy}=0$，则所有 $y_i$ 必须相同。

（2）任意满足 $\sum_{i=1}^{n}(y_i-a-bx_i)=0$ 的实参数对 $(a,b)$，都是最小二乘解。

（3）若某实参数对同时满足两条正规方程，则它是唯一最小二乘解。

答案统一置于讲义末尾。
\end{exercise}
\section{结果检验、退化情形与工程解释}\label{sec:S03}
% UNIT S03.U01
第\ref{sec:S02}节给出了带截距直线的求解方法。本节进一步回答三个问题：怎样利用残差检验计算结果？当自变量观测值全相同时还能确定什么？把拟合直线用于工程预测时，哪些结论不能仅由平方和最小推出？以下仍只使用有限求和、二元一次方程和配方法。\par
% UNIT S03.U02
\begin{definition}
设 $n\geq1$ 为整数，观测数据为 $(x_i,y_i)\in\mathbb{R}^2$，$i=1,\ldots,n$。选定拟合函数 $f$ 后，令 $\hat{y}_i=f(x_i)$、$r_i=y_i-\hat{y}_i$。残差表是按观测编号列出 $x_i,y_i,\hat{y}_i,r_i$ 的表格，也可以增加 $r_i^2$ 一列，以便求出残差平方和。

插值要求函数在每个观测位置都准确命中观测值，即 $f(x_i)=y_i$ 对所有 $i$ 成立。拟合则是在指定函数模型内按某种整体准则选择参数，不要求逐点命中。本讲义采用的准则是最小化 $\sum_{i=1}^{n}r_i^2$。若同一个输入对应不同的观测输出，就不存在同时命中这些记录的单值插值函数；最小二乘拟合仍有意义。

记 $x_{\min}=\min\{x_1,\ldots,x_n\}$、$x_{\max}=\max\{x_1,\ldots,x_n\}$。将拟合模型用于输入 $x_0\in[x_{\min},x_{\max}]$ 的输出估计，称为区间内预测；用于 $x_0<x_{\min}$ 或 $x_0>x_{\max}$ 的估计，称为外推。区间内预测并不等于要求逐点命中的插值。

预测所用的 $\hat{y}(x_0)=f(x_0)$ 是模型计算值。即使 $x_0$ 位于已有观测位置，它也不一定等于该处已有的观测值，更不能保证等于以后在同一输入下测得的实际输出。
\end{definition}
% UNIT S03.U03
\begin{theorem}
设 $n\geq1$ 为整数，$(x_i,y_i)\in\mathbb{R}^2$，$i=1,\ldots,n$，且 $\bar{x}=\frac1n\sum_{i=1}^{n}x_i$、$\bar{y}=\frac1n\sum_{i=1}^{n}y_i$。若 $(a_*,b_*)$ 是目标函数
\[
Q(a,b)=\sum_{i=1}^{n}(y_i-a-bx_i)^2,\qquad (a,b)\in\mathbb{R}^2,
\]
的一个全局最小点，则其残差 $r_i=y_i-a_*-b_*x_i$ 满足
\[
\sum_{i=1}^{n}r_i=0,\qquad a_*+b_*\bar{x}=\bar{y}.
\]
因此最小二乘直线经过数据均值点 $(\bar{x},\bar{y})$。这些结论不要求 $S_{xx}>0$，但要求这里的模型包含可自由选择的截距，且使用上述未加权的残差平方和。
\end{theorem}
% UNIT S03.U04
\begin{proof}
固定最优参数中的斜率 $b_*$。由第\ref{sec:S01}节的常数拟合结论，数据 $y_i-b_*x_i$ 的唯一最优常数是其平均数 $\bar{y}-b_*\bar{x}$。如果 $a_*$ 不等于这个数，仅改变截距就会严格减小平方和，与全局最优性矛盾。因此
\[
a_* =\bar{y}-b_*\bar{x}.
\]
这就是第\ref{sec:S02}节第一条正规方程除以 $n$ 后的形式，也直接说明拟合直线经过均值点。再按残差定义求和，得到
\[
\sum_{i=1}^{n}r_i
=\sum_{i=1}^{n}y_i-na_*-b_*\sum_{i=1}^{n}x_i
=n(\bar{y}-a_*-b_*\bar{x})=0.
\]
整个论证没有除以 $S_{xx}$，因而也适用于自变量观测值全相同的情形。
\end{proof}
% UNIT S03.U05
残差和为零与通过均值点实际上是同一个等式的两种表达，不能把它们当作两项独立的最优性条件。残差和为零也不表示每个残差都为零。若 $S_{xx}>0$，还可以检验第二条正规方程，即 $\sum_{i=1}^{n}x_i r_i=0$；由第\ref{sec:S02}节的结论，同时满足两条正规方程的参数对才是唯一最小二乘解。

这些代数条件用于检验是否求对了指定模型的最优参数，并不证明模型适合工程对象。若按输入从小到大观察残差，发现它们呈现有规律的变化，可将其作为重新考虑直线模型的线索；但少量残差本身不足以确定变化原因。

截距的前提不能省略。例如，仅允许过原点的模型 $\hat{y}=bx$ 时，对数据 $(0,1)$、$(1,1)$ 有
\[
Q(b)=1+(1-b)^2,\qquad b\in\mathbb{R}.
\]
唯一最优斜率为 $b=1$，残差为 $1,0$，其和不为零；直线 $\hat{y}=x$ 也不经过数据均值点 $(1/2,1)$。\par
% UNIT S03.U06
\begin{example}
把第\ref{sec:S02}节的数据 $(0,1)$、$(1,2)$、$(2,2)$ 看成某传感器的三组输入输出记录，输入和输出均按各自固定单位记数。已求得最小二乘直线 $\hat{y}=7/6+x/2$。列出残差表，核验残差和与均值点性质；估计输入 $1.5$ 时的输出，并讨论把模型用于输入 $10$ 时的局限。
\end{example}
% UNIT S03.U07
\par\noindent\textbf{解：}对实数输入 $x$，拟合模型为 $\hat{y}(x)=7/6+x/2$。在已有观测处，残差表为
\[
\begin{array}{c|r|r|r|r|r}
i&x_i&y_i&\hat{y}_i&r_i&r_i^2\\ \hline
1&0&1&7/6&-1/6&1/36\\
2&1&2&5/3&1/3&1/9\\
3&2&2&13/6&-1/6&1/36
\end{array}
\]
因此
\[
\sum_{i=1}^{3}r_i=-\frac16+\frac13-\frac16=0,
\qquad Q(a_*,b_*)=\frac16.
\]
负残差表示模型高估了该次观测输出，正残差表示模型低估了该次观测输出。这里模型没有命中任何一个观测点，却仍使总平方和最小。

均值点为 $(\bar{x},\bar{y})=(1,5/3)$，且
\[
\hat{y}(1)=\frac76+\frac12=\frac53=\bar{y}.
\]
第二条残差关系也成立：$\sum_{i=1}^{3}x_i r_i=0+1/3-2/6=0$。

输入 $1.5=3/2$ 位于观测输入区间 $[0,2]$ 内，模型给出的区间内预测为
\[
\hat{y}\left(\frac32\right)=\frac76+\frac12\cdot\frac32
=\frac{23}{12}.
\]
这是按输出单位记数的模型估计，不是对未来实际读数的保证。

若输入为 $10$，代入同一直线得到
\[
\hat{y}(10)=\frac76+\frac12\cdot10=\frac{37}{6}.
\]
但 $10$ 远在 $[0,2]$ 之外，这属于外推。传感器在较大输入处可能不再近似线性，例如可能出现饱和；现有三组记录不能排除这些情况。平方和最小只描述模型对已有记录的整体误差，既不保证外推可靠，也不能单凭这次拟合赋予模型因果含义。即使在观测区间内使用，也应结合传感器工作条件与新的测量记录判断适用性。\par
% UNIT S03.U08
\begin{theorem}
设 $n\geq1$ 为整数，$x_0,y_1,\ldots,y_n\in\mathbb{R}$，且所有自变量观测值均为 $x_i=x_0$。令 $\bar{y}=\frac1n\sum_{i=1}^{n}y_i$。对带截距直线模型，有
\[
Q(a,b)=\sum_{i=1}^{n}(y_i-\bar{y})^2+n(a+bx_0-\bar{y})^2,
\qquad (a,b)\in\mathbb{R}^2.
\]
全部最优参数对恰为
\[
a=\bar{y}-bx_0,\qquad b\in\mathbb{R}.
\]
因此最优参数对不唯一，最小平方和为 $\sum_{i=1}^{n}(y_i-\bar{y})^2$，但所有最优直线在共同输入 $x_0$ 处的拟合值唯一，等于 $\bar{y}$。对每一个固定的实数 $x\ne x_0$，这些最优直线在 $x$ 处的拟合值不唯一。
\end{theorem}
% UNIT S03.U09
\begin{proof}
因为每个 $x_i=x_0$，数据处的拟合值都是同一个数 $c=a+bx_0$，所以
\[
Q(a,b)=\sum_{i=1}^{n}(y_i-c)^2.
\]
任意实数 $c$ 都可通过取 $a=c,b=0$ 实现。由第\ref{sec:S01}节的常数拟合平方和分解，立刻得到本定理的分解式。由于 $n>0$，该式在且仅在 $a+bx_0=\bar{y}$ 时取到最小值。

对每个 $b\in\mathbb{R}$，取 $a=\bar{y}-bx_0$ 就得到一个最优参数对；不同的 $b$ 给出不同的参数对，故共有无穷多个。它们在 $x_0$ 处的拟合值都是 $\bar{y}$。而在任意固定的 $x\ne x_0$ 处，
\[
\hat{y}(x)=\bar{y}-bx_0+bx=\bar{y}+b(x-x_0).
\]
因为 $x-x_0\ne0$，改变 $b$ 会改变该拟合值。事实上，对任意给定的 $t\in\mathbb{R}$，取 $b=(t-\bar{y})/(x-x_0)$ 就有 $\hat{y}(x)=t$。因此这些数据并未唯一确定其他输入位置的拟合输出。
\end{proof}
% UNIT S03.U10
\begin{example}
对数据点 $(2,1)$、$(2,3)$、$(2,5)$ 作带截距的直线最小二乘拟合。求全部最优参数对、残差和最小平方和。能否唯一确定输入 $2$ 与输入 $3$ 时的拟合输出？
\end{example}
% UNIT S03.U11
\par\noindent\textbf{解：}三个自变量观测值均为 $2$，所以 $\bar{x}=2$、$S_{xx}=0$，不能使用 $b_*=S_{xy}/S_{xx}$ 的除法公式。响应平均数为 $\bar{y}=(1+3+5)/3=3$。对所有 $a,b\in\mathbb{R}$，
\[
\begin{aligned}
Q(a,b)&=(1-a-2b)^2+(3-a-2b)^2+(5-a-2b)^2\\
&=8+3(a+2b-3)^2.
\end{aligned}
\]
因此全部最优参数满足
\[
a+2b=3,
\quad\text{即}\quad (a,b)=(3-2b,b),\qquad b\in\mathbb{R}.
\]
每组最优参数在输入 $2$ 处都给出拟合值 $3$，对应残差依次为 $-2,0,2$，最小平方和为 $4+0+4=8$。

例如 $\hat{y}=3$ 与 $\hat{y}=1+x$ 都是最优直线。它们在输入 $3$ 处分别给出 $3$ 与 $4$。一般地，最优直线在实数输入 $x$ 处的值为
\[
\hat{y}(x)=3+b(x-2),\qquad b\in\mathbb{R}.
\]
所以输入 $2$ 时的拟合输出唯一，输入 $3$ 时的拟合输出却随斜率改变。本例并不是没有最小二乘解，而是数据只约束了直线在一个输入位置的高度，无法确定斜率。\par
% UNIT S03.U12
处理一组数据时，应区分三个层次：先确认参数是否使指定目标函数最小，再看残差反映了怎样的逐点误差，最后判断模型是否适合预定用途。参数唯一性也必须单独考虑：$S_{xx}>0$ 时带截距的最小二乘直线唯一；$S_{xx}=0$ 时最优参数不唯一，但观测输入处的拟合值仍唯一。

工程报告宜同时交代观测输入范围、所用模型、参数、残差平方和以及预测是否属于外推。拟合公式不是物理规律的自动证明，残差和为零也不是误差已经消失。\par
% UNIT S03.U13
\begin{exercise}
\label{mathagent:exercise:792316d6ae426af3441a87022eb342411ffcb273e3da7957f4be83fb59518d0a}
某装置的三组输入输出记录为 $(0,1)$、$(1,1)$、$(2,3)$，其带截距最小二乘直线为 $\hat{y}=2/3+x$。列出拟合值和残差，核验残差和与均值点性质；求输入 $3/2$ 与输入 $4$ 时的模型输出，分别说明属于区间内预测还是外推，并解释这些数值是否必然等于未来实测输出。答案统一置于讲义末尾。
\end{exercise}
% UNIT S03.U14
\begin{exercise}
\label{mathagent:exercise:2a86525c0362bd6004ffb20b570bb82a80598d07bb305d4bd130c7cd7cc87533}
对数据点 $(1,2)$、$(1,4)$ 作带截距直线最小二乘拟合。求全部最优参数、最小平方和以及输入 $1$ 处的拟合值，并给出两条在输入 $0$ 处输出不同的最优直线。由此判断：“只要最优参数不唯一，已有观测位置的拟合值就一定不唯一。”答案统一置于讲义末尾。
\end{exercise}
\section*{练习答案}
% UNIT S01.U10.answer
\par\noindent\textbf{\hyperref[mathagent:exercise:79e91357ee05b8b9aa4c657d553157f1ef411624d282c3f2337dd0e01453ce2d]{练习 \ref*{mathagent:exercise:79e91357ee05b8b9aa4c657d553157f1ef411624d282c3f2337dd0e01453ce2d} 的解答：}}目标函数为
\[
Q(c)=(1-c)^2+(2-c)^2+(4-c)^2+(5-c)^2,\qquad c\in\mathbb{R}.
\]
四次测量值的平均数为
\[
\bar{y}=\frac{1+2+4+5}{4}=3.
\]
由常数拟合的平方和分解，
\[
\begin{aligned}
Q(c)&=(1-3)^2+(2-3)^2+(4-3)^2+(5-3)^2+4(c-3)^2\\
&=10+4(c-3)^2.
\end{aligned}
\]
因此唯一最优常数为 $c=3$，对应的残差依次为
\[
r_1=-2,\qquad r_2=-1,\qquad r_3=1,\qquad r_4=2,
\]
最小平方和为 $Q(3)=4+1+1+4=10$。

当 $c=2$ 时，
\[
Q(2)=(-1)^2+0^2+2^2+3^2=14.
\]
平方和分解给出
\[
Q(2)-Q(3)=4(2-3)^2=4.
\]
也就是说，候选常数 $2$ 偏离平均数 $3$ 一个单位，使残差平方和比最小值增加了 $4$。\par
% UNIT S01.U11.answer
\par\noindent\textbf{\hyperref[mathagent:exercise:9e7dbf7ef17df1506f236bbf1ad6ff8c14ff13065df89e12b796d18475228be7]{练习 \ref*{mathagent:exercise:9e7dbf7ef17df1506f236bbf1ad6ff8c14ff13065df89e12b796d18475228be7} 的解答：}}（1）正确。由残差定义与平均数定义，
\[
\sum_{i=1}^{n}r_i=\sum_{i=1}^{n}y_i-nc=n(\bar{y}-c).
\]
因为 $n>0$，残差之和为零必有 $c=\bar{y}$。由常数拟合定理，$\bar{y}$ 正是该模型的唯一最优参数。

（2）错误。残差之和为零可以由正负残差抵消产生。例如取 $n=2$、$y_1=1$、$y_2=3$ 及 $c=2$，则
\[
r_1=-1,\qquad r_2=1,\qquad r_1+r_2=0,
\]
但两个残差均不为零。因此题设不能推出每个残差都为零。

（3）正确。当 $n=1$ 时，
\[
Q(c)=(y_1-c)^2\geq0.
\]
等号当且仅当 $c=y_1$ 成立，所以唯一最优参数为 $c=y_1=\bar{y}$，最小平方和为零。\par
% UNIT S02.U11.answer
\par\noindent\textbf{\hyperref[mathagent:exercise:332ae2e99c79cf97370ba4732967ff23f2a2c4a4b36b98eda8e910755bf43b8b]{练习 \ref*{mathagent:exercise:332ae2e99c79cf97370ba4732967ff23f2a2c4a4b36b98eda8e910755bf43b8b} 的解答：}}均值为 $\bar{x}=1$、$\bar{y}=5/3$。中心化数据依次为
\[
(x_i-\bar{x})_{i=1}^{3}=(-1,0,1),\qquad
(y_i-\bar{y})_{i=1}^{3}=(-2/3,-2/3,4/3).
\]
所以
\[
S_{xx}=1+0+1=2,\qquad S_{xy}=\frac23+0+\frac43=2.
\]
由此 $b_*=1$、$a_*=2/3$，最小二乘直线为 $\hat{y}=2/3+x$。拟合值依次为 $2/3,5/3,8/3$，残差依次为 $1/3,-2/3,1/3$，最小平方和为
\[
Q(a_*,b_*)=\frac19+\frac49+\frac19=\frac23.
\]
本题的正规方程是 $3a+3b=5$、$3a+5b=7$；代入 $a=2/3,b=1$，左侧分别为 $2+3=5$ 与 $2+5=7$，均与右侧相等。\par
% UNIT S02.U12.answer
\par\noindent\textbf{\hyperref[mathagent:exercise:b4c9c9a1fb82bcef052fdbf6e1539831ed88cee016dbcbf83a50e5e8d2311be0]{练习 \ref*{mathagent:exercise:b4c9c9a1fb82bcef052fdbf6e1539831ed88cee016dbcbf83a50e5e8d2311be0} 的解答：}}（1）错误。例如数据点 $(-1,1)$、$(0,0)$、$(1,1)$ 满足 $\bar{x}=0$、$\bar{y}=2/3$、$S_{xx}=2>0$，且
\[
S_{xy}=(-1)\frac13+0\left(-\frac23\right)+1\frac13=0,
\]
但三个响应值不全相同。此时最优直线是 $\hat{y}=2/3$。

（2）错误。残差和为零仅等价于 $a+b\bar{x}=\bar{y}$，即截距对于当前斜率已最优，却没有保证斜率也最优。例如对本节例题中的数据 $(0,1)$、$(1,2)$、$(2,2)$，取 $a=5/3,b=0$，残差和为零，但平方和为 $2/3$，大于最小值 $1/6$。

（3）正确。在 $S_{xx}>0$ 的前提下，消元使正规方程等价于 $a+b\bar{x}=\bar{y}$ 与 $bS_{xx}=S_{xy}$，故其唯一解为 $b_*=S_{xy}/S_{xx}$、$a_*=\bar{y}-b_*\bar{x}$。由全局平方和分解，这正是唯一最小二乘解。\par
% UNIT S03.U13.answer
\par\noindent\textbf{\hyperref[mathagent:exercise:792316d6ae426af3441a87022eb342411ffcb273e3da7957f4be83fb59518d0a]{练习 \ref*{mathagent:exercise:792316d6ae426af3441a87022eb342411ffcb273e3da7957f4be83fb59518d0a} 的解答：}}拟合值依次为 $2/3,5/3,8/3$，残差依次为 $1/3,-2/3,1/3$，所以残差和为零。均值为 $\bar{x}=1$、$\bar{y}=5/3$，且 $\hat{y}(1)=2/3+1=5/3$，故直线经过均值点。

两个模型输出分别为
\[
\hat{y}\left(\frac32\right)=\frac23+\frac32=\frac{13}{6},\qquad
\hat{y}(4)=\frac23+4=\frac{14}{3}.
\]
观测输入区间为 $[0,2]$，故前者是区间内预测，后者是外推。二者都是模型计算值，不必等于未来实测输出；尤其对输入 $4$，现有数据不能保证区间外仍遵循同一近似线性关系。\par
% UNIT S03.U14.answer
\par\noindent\textbf{\hyperref[mathagent:exercise:2a86525c0362bd6004ffb20b570bb82a80598d07bb305d4bd130c7cd7cc87533]{练习 \ref*{mathagent:exercise:2a86525c0362bd6004ffb20b570bb82a80598d07bb305d4bd130c7cd7cc87533} 的解答：}}共同输入为 $1$，响应平均数为 $3$。对所有 $a,b\in\mathbb{R}$，
\[
Q(a,b)=(2-a-b)^2+(4-a-b)^2=2+2(a+b-3)^2.
\]
所以全部最优参数为 $a=3-b$、$b\in\mathbb{R}$，最小平方和为 $2$。所有最优直线在输入 $1$ 处的拟合值都是 $3$，残差为 $-1,1$。

例如 $\hat{y}=3$ 与 $\hat{y}=2+x$ 都最优，但在输入 $0$ 处分别给出 $3$ 与 $2$。题中判断错误：参数不唯一并不必然导致观测位置的拟合值不唯一，本题就是反例。\par
\end{document}
